\NewDocumentCommand\AxiomGR{}{GR}
\NewDocumentCommand\PrintAxiomGR{}{
  \begin{axiom}[\AxiomGR]
    Every proposition is definable: $\fa p\Omega \ex{n}{\ovob} p = \mbb p_n$.
  \end{axiom}
}


\chapter{Quasi-coherence for the propositional universe}\label{ch:sqcforomega}

In this chapter we explore the possibility of a general theory of quasi-coherence that characterises the propositional universe $\Omega$ in the classifying topos of a restrictive class of geometric theories, which we call \emph{countably generated} (cf.~\zcref{def:countablyseparated}). We believe this is an initial step towards a general theory of quasi-coherence that works for arbitrary geometric theories, but we restrict ourselves to the special case this time.

We comment that interestingly, the \emph{Nullstellensatz} result holds \emph{constructively}, while the full \gls{QC} principle we obtain for the propositional universe $\Omega$ does \emph{not} --- we will discuss a counterexample at~\zcref{rem:roleofexcludedmiddle}. In particular, some of the results in this chapter requires \gls{LEM} in the meta-theory, in which cases we will indicate explicitly. 

\section{Nullstellensatz for countably generated geometric theories}

Recall the notion of a countable set in~\zcref{exm:countablesize}. We say a set $A$ is \emph{enumerable} if there exists a surjection $\N \surj 1 + A$. All countable sets are thus enumerable.

Let $\T$ be a geometric theory. By a \emph{primitive} $\T$-formula we mean one of the form 
\[ \exists\ov{\ms y}.\,\bigwedge_{i\in n}\varphi_i(\ov{\ms x},\ov{\ms y}), \]
where each $\varphi_i$ is atomic; notice that we count $\bot$ as an atomic formula. The set of primitive formulas in $\T$ in context $\ov{\ms x}$ will be denoted as $\form^p_\T(\ov{\ms x})$. We will use $\form^\N_\T(\ov{\ms x})$ to denote the set of geometric formulas in $\T$ that are $\N$-disjunctions of formulas in $\form^p_\T(\ov{\ms x})$. Equivalently, an element in $\form^\N_\T(\ovs x)$ corresponds to an infinite sequence $\varphi \in \form^p_\T(\ovs x)^\N$, where we view $\varphi = \bigvee_{i\in\N}\varphi_i$ as an infinite disjunction of primitive formulas. 

\begin{remark}[Normal forms of geometric formulas]\label{rem:normalformofgeometric}
  It is well-known that geometric formulas has a normal form, where every geometric formula can be written as a disjunction of primitive formulas. This can be achieved by continually commute existential quantifiers past (infinite) disjunctions, and take the disjunctive normal forms of quantifier free formulas.
\end{remark}

\begin{definition}[Countably generated geometric theories]\label{def:countablyseparated}
  We say a geometric theory $\T$ is \emph{countably generated} if the following holds:
  \begin{itemize}
    \item it has a countable signature, i.e.\ its set of sorts, function symbols, and relation symbols are all countable;
    \item every axiom of $\T$ is of the form $\varphi \vdash_{\ovs x} \psi$, where $\varphi$ is atomic and $\psi \in \form^\N_\T(\ovs x)$.\footnote{Note that the requirement of $\varphi$ being atomic can be weakened by it being in $\form^\N_\T(\ovs x)$ as well, since the appearance of disjunctions and existential quantifiers on the left can always be reduced in geometric logic.}
  \end{itemize}
\end{definition}

\begin{remark}
  By~\zcref{rem:normalformofgeometric}, if we assume \gls{LEM}, then every geometric formula in a countably generated theory $\T$ will be equivalent to one in $\form^\N_\T$.
\end{remark}

\begin{example}[The theory of filters on $\N$]\label{exa:toposoftrees}
  Recall that $\N$ under the canonical order is a linear poset. The theory $\mbb F_{\N}$ of a \emph{filter on $\N$} is a propositional theory defined as follows:
  \begin{itemize}
    \item it has countably many propositional variables $\ms p_0,\ms p_1,\cdots$, where each $\ms p_i$ intuitively states that $i$ belongs to the filter;
    \item for any $i \le j$, a propositional axiom $\ms p_i \vdash \ms p_j$, stating the filter is upward closed;
    \item a single axiom $\top \vdash \bigvee_{i\in n}\ms p_i$, stating the filter is inhabited.
  \end{itemize}
  By construction, this theory $\mbb F_{\N}$ is countably generated. Its classifying topos is the presheaf topos $\psh(\N)$ on the poset $\N$, and is often called the \emph{topos of trees}.
\end{example}

\begin{example}[The theory of local, locally finite, algebraically closed Heyting algebras]
  Following~\citet[\S 7.3]{ghilardi2013sheaves}, the theory $\mbb K$ extends the algebraic theory $\mbb H$ of Heyting algebras with additional axioms. For an finite conjunction of equations $\varphi(\ov{\ms x})$ in the language of Heyting algebras, we say it is \emph{finite} if the corresponding finitely presented Heyting algebra $H[\ov{\ms x}]/\varphi(\ov{\ms x})$ is \emph{finite}.\footnote{By the presentation of finitely presented Heyting algebras, whether a formula $\varphi(\ov{\ms x})$ is finite is \emph{decidable}.} The set of finite formulas with $n$ free variables will be denoted as $F_n$. $\mbb K$ then has the following axioms:
  \begin{itemize}
    \item $\top \vdash \bigvee\scomp{\varphi(\ov{\ms x})}{\varphi(\ov{\ms x})\in F_n}$, for any $n\in\N$;
    \item for any $\varphi(\ov{\ms x})$ and $\psi(\ov{\ms y})$ such that the induced map $H[\ov{\ms x}]/\varphi(\ov{\ms x}) \to H[\ov{\ms x},\ov{\ms y}]/\psi(\ov{\ms x},\ov{\ms y})$ is an injection, an axiom $\varphi(\ov{\ms x}) \vdash_{\ov{\ms x}} \exists\ov{\ms y}.\,\psi(\ov{\ms x},\ov{\ms y})$;
    \item locality axioms $\top \vdash 0 \neq 1$ and $\ms x \vee \ms y = 1 \vdash_{\ms x,\ms y} \ms x = 1 \vee \ms y = 1$.
  \end{itemize}
  By construction, $\mbb K$ is countably generated. Its classifying topos is constructed in~\citet{ghilardi2013sheaves}, and is called the K-topos in~\citet{ye2026coexact}. 
\end{example}

We end this section by mentioning the \emph{Nullstellensatz} specifically for a countably generated geometric theory $\T$. Let $\Set[\T]$ be its classifying topos with the generic model $V_\T$. For any geometric formula $\varphi(\ovs x)$, let $\ass{\varphi(\ovs x)}$ denote its interpretation in $V_\T$.

Similar to the construction of~\zcref{def:theoryofdiagram}, we have an internal theory of diagrams $\VT$ in $\Set[\T]$. Thus, we can similarly define the internal sets $\form^p_{V_\T\ltimes\T}$ and $\form^\N_{V_\T\ltimes\T}$ of primitive formulas and their countable disjunctions in $\VT$, respectively. In particular, by definition $\form^\N_{V_\T\ltimes\T}(\ovs x) \cong \form^p_{V_\T \ltimes \T}(\ovs x)^\N$ internally in $\Set[\T]$. 

\begin{theorem}[Nullstellensatz]\label{thm:Nullstellensatz}
  The following holds internally in $\Set[\T]$,
  \[ \fa{\varphi,\psi}{\form^\N_{V_\T \ltimes \T}(\ovs x)} V_\T \models (\varphi \vdash_{\ovs x} \psi) \eff \VT \vdash (\varphi \vdash_{\ovs x} \psi). \]
\end{theorem}
\begin{proof}
  The proof of~\citet[Thm.~3.7]{blechschmidt_general_2020} holds here, especially because we have restricted the internal geometric formulas to those having disjunctions indexed by $\N$, which is locally constant in the sense of \emph{loc.\ cit.}
\end{proof}


\section{Axiomatising the propositional universe}

\zcref{thm:Nullstellensatz} is a \emph{conservativity} result, which states that for sequents in $\form^\N_{V_\T\ltimes\T}$ internally in $\Set[\T]$, it is provable if and only if it is true at the generic $\T$-model $V_\T$. In particular, this implies that we have an \emph{embedding} internally in $\Set[\T]$, that
\[ \form^\N_{V_\T\ltimes\T}(\ovs x)/\vneg \hook \mc P\ass{\top(\ovs x)}, \]
where $\mc P(-) \coloneq \Omega^-$ denotes the internal power construction, and $\sim$ is the equivalence relation generated by $\VT$-provability. 

Our goal in this section is to show that the above map is an \emph{isomorphism}. However, our chosen formulation only holds when we assume \gls{LEM} in the external meta-theory, which we will freely do so from now on and indicate its use explicitly. We will comment on the use of \gls{LEM} in~\zcref{rem:differentdefinability}. We first observe that we have an external definability result:

\begin{lemma}[External definability of subobjects of the generic model]\label{lem:externaldef}
  Any subobject $S \hook \ass{\top(\ovs x)}$ is defined by an external geometric formula $\beta_S(\ovs x)$, where concretely we have 
  \[ \beta_S(\ovs x) \coloneq \bigvee\scomp{\gamma(\ovs x)}{\gamma(\ovs x) \in \form^p_\T(\ovs x) \conjt \ass{\rho(\ovs x)} \le S}. \]
\end{lemma}
\begin{proof}
  As mentioned in~\zcref{rem:normalformofgeometric}, every geometric formula is a disjunction of formulas in $\form^p_\T(\ovs x)$. This then follows from the definability result in~\citet[Thm.~2.4]{caramello2012universal}. 
\end{proof}

\begin{remark}\label{rem:subcomlem}
  Assuming \gls{LEM}, every $\beta_S(\ovs x)$ indeed belongs to $\form^\N_\T(\ovs x)$, since by construction $\form^p_\T(\ovs x)$ itself will be countable since $\T$ is countably generated.
\end{remark}

\begin{theorem}[Definability of the propositional universe for countably generated theories]\label{thm:omegadefinable}
  Assuming \gls{LEM} in the meta-theory, for any countably generated theory $\T$, the map obtained by interpreting a $\VT$-formula in the generic $\T$-model $V_\T$ induces an isomorphism
  \[ \form^\N_{V_\T\ltimes\T}(\ovs x)/\vneg \cong \mc P\ass{\top(\ovs x)}. \]
\end{theorem}
\begin{proof}
  We have mentioned that by~\zcref{thm:Nullstellensatz}, this map is an embedding, thus it suffices to show it is a surjection. By Kripke-Joyal semantics, consider any $A \coloneq \alpha(\ovs y) \in \mb C[\Sigma]$, and consider an object $S \in \mc P\ass{\top(\ovs x)}(A) \cong \sub(\yon_A \times \ass{\top(\ovs x)})$. Now by construction, canonically we have that 
  \[ a \colon \yon_A \cong \ass{\alpha(\ovs y)} \hook \ass{\top(\ovs y)}, \] 
  which we may view as the generic $A$-element of sorts $\ovs y$. Thus, we may then view $S$ as a subobject of $\ass{\top(\ovs y)} \times \ass{\top(\ovs x)} \cong \ass{\top(\ovs y,\ovs x)}$. By~\zcref{lem:externaldef} and~\zcref{rem:subcomlem}, assuming \gls{LEM} we can find an external formula $\beta_S \in \form^\N_\T(\ovs y,\ovs x)$ such that $\ass{\beta_S} = S$. By substituting the generic term $a$ into $\beta_S$, we obtain an internal formula $\beta_S(a,\ovs x) \in \form^\N_{V_\T\ltimes\T}(\ovs x)$, such that by construction $\ass{\beta_S(a,\ovs x)} \cong S$. Hence, this map is also surjective.
\end{proof}

\begin{remark}[The difference with~\citet{blechschmidt_general_2020}]\label{rem:differentdefinability}
  It is already mentioned in~\citet[Rem.~3.8]{blechschmidt_general_2020} that the Nullstellensatz result as stated in~\zcref{thm:Nullstellensatz} cannot hold for internal geometric formulas with disjunctions indexed by arbitrary internal sets in $\Set[\T]$, thus \emph{loc.\ cit.}\ tries to fix this by introducing the notion of \emph{locally constant} indexing sets. However, to our best knowledge this notion, though externally makes sense, is \emph{not internally definable}, at least not obviously so. Thus, we have restricted Nullstellensatz to the case for formulas in $\form^\N_{V_\T\ltimes\T}$, whose disjunctions are indexed by $\N$, which is indeed a globally \emph{constant} object in any sheaf topos.

  However, the definability result we have obtained in~\zcref{thm:omegadefinable} becomes \emph{stronger} than that of Thm.~5.2 in \emph{loc.\ cit.}, which states that $\mc P\ass{\top(\ovs x)}$ can be internally defined by geometric formulas in $\VT$ with locally constant disjunctions. Instead, our definability result only uses \emph{countably indexed disjunctions}, which is a stronger statement \emph{a priori}. However, this restriction is related to why we need to assume \gls{LEM} in the meta-theory when establishing~\zcref{thm:omegadefinable}, which we comment in~\zcref{rem:roleofexcludedmiddle}.
\end{remark}

\begin{remark}[The role of excluded middle]\label{rem:roleofexcludedmiddle}
  It is very interesting that we need \gls{LEM} in the external meta-theory for~\zcref{thm:omegadefinable} to hold. By inspecting its proof, we have mainly used \gls{LEM} to transform an arbitrary subset of $\N$ into a countable subset, thus can be lifted to an internal formula in $\form^\N_{V_\T\ltimes\T}$ in $\Set[\T]$. As a comparison, the corresponding internal definability result in~\citet{blechschmidt_general_2020}, using all locally constant sheaves as indexing sets for disjunctions, though not fully formal, is obtained constructively.

  However, the simplest example would indicate why \gls{LEM} is necessary for our stronger definability result. Take the simplest geometric theory $\mbb 1$ of a \emph{singleton}, which has a single sort, no relation or function symbols, and two axioms 
  \[ \top \vdash \exists \ms x.\,\ms x \quad \top \vdash_{\ms x,\ms y} \ms x = \ms y, \]
  stating that there is exactly one object. By construction, its classifying topos is the terminal topos $\Set[\mbb 1] \simeq \Set$, since it classifies the terminal object in a topos. The diagram theory $V_{\mbb 1} \!\ltimes\! \mbb 1$ this time is simply $\mbb 1$ extended by a single constant $\ms c$, denoting the unique element. It is then easy to observe that any primitive sentence in $V_{\mbb 1} \!\ltimes\! \mbb 1$ is provably equivalent to either $\top$ or $\bot$. If the definability result in~\zcref{thm:omegadefinable} holds, then we would have a diagram 
  \[
  \begin{tikzcd}
    \form^\N_{V_{\mbb 1} \ltimes \mbb 1}(\emp) \cong \form^p_{V_{\mbb 1}\ltimes\mbb 1}(\emp)^\N \ar[r, two heads] \ar[dr, two heads] & 2^\N \ar[d, dashed, two heads, "q"] \\
    & \Omega \cong \pp
  \end{tikzcd}
  \]
  where by construction, $q$ is the map 
  \[ q(f) = \ex n{\N} f(n) = 1. \]
  In particular, this implies that all propositions in $\Set \simeq \Set[\mbb 1]$ is \emph{semi-decidable} (cf.~\zcref{rem:dominanceofsemidecidable}), which of course is \emph{not} necessarily true without \gls{LEM}!
\end{remark}

In the remaining of this chapter we mention two applications of \gls{QC} for the propositional universe as established in~\zcref{thm:omegadefinable}. We haven't put this into~\zcref{prt:appl} of this thesis because they are relatively small, instead of developing a whole field of synthetic mathematics. However, we do find them useful and interesting, thus they are included here.

\section{Application I: the topos of trees}

Recall from~\zcref{exa:toposoftrees} that we have a countably generated geometric theory $\mbb F_\N$ of filters on the linear order $\N$, whose classifying topos is given by the topos of trees 
\[ \Set[\mbb F_\N] \simeq \psh(\N). \]
This topos is used in~\citet{birkedal_first_2012} to model \emph{synthetic guarded domain theory (\gls{SGDT})}, and one crucial fact is that there is a \emph{modality} on the universe of propositions in $\psh(\N)$ which satisfies \emph{L\"ob induction}. We will see that these structures and their properties can be \emph{defined} and \emph{proved} completely by the internal definability result of \gls{QC} for the coherent theory $\mbb F_\N$.

Recall that $\set{\ms p_n}_{n\in\N}$ is the family of propositional predicates in $\mbb F_{\N}$. We also write $\mbb p_n$ for the interpretation $\ass{\ms p_n}$ internally in $\psh(\N)$, and we write $\mbb p$ as the generic model. Since a presheaf topos is \emph{locally connected} (cf.~\citet[C3.3]{johnstone2002sketches}), by assuming \gls{LEM} in the meta-theory, internally the \emph{Limited Principle of Omniscience (\gls{LPO})} holds,
\[ \psh(\N) \models \fa f{2^\N} (\ex n\N fn = 0) \vee (\fa n\N fn = 1). \]

\begin{proposition}\label{prop:qcforFN}
  The following holds internally in $\psh(\N)$ by \gls{QC} for the propositional universe $\Omega$ (under \gls{LPO}):
  \[ \fa p\Omega p = 0 \vee \ex n{\N} p = \mbb p_n. \]
\end{proposition}
\begin{proof}
  Notice that since $\mbb F_\N$ is a propositional theory having no sorts, there are no additional constants added in $\mbb F_\N$. This shows that primitive formulas of $\mbb p \ltimes \mbb F_\N$ coincides with that in $\mbb F_\N$, which are simply finite conjunctions and disjunctions of $\set{\ms p_i}_{i\in\N} \cup \set{\bot,\top}$. In fact, up to provability in $\mbb p \ltimes \mbb F_\N$, we can choose the primitive formulas in $\mbb p\ltimes\mbb F_\N$ to be $P \coloneq \set{\ms p_i}_{i\in\N} \cup \set{\bot}$, since by the axiomatisation in $\mbb F_\N$, $\top$ can be defined by a countable disjunction $\bigvee_{i\in\N}\mbb p_i$. Hence, by~\zcref{thm:omegadefinable} we have
  \[ \fa p{\Omega} \ex {f}{P^\N} p = \bigvee_{i\in\N} \ass{f_i}. \]
  Notice that $P$ is naturally a linear order isomorphic to $\N$. Under \gls{LPO}, we can decide whether a sequence $f\in P^\N$ is bounded or not. In the former case, there exists a maximal element in the image of $f$, which is either $\bot$ or some $\ms p_n$. Hence, $\bigvee_{i\in\N}\ass{f_i} = \mbb p_n$. In the latter case, $p = \bigvee_{i\in\N}\ass{f_i} = \bigvee_{i\in\N}\mbb p_n$. Hence, it also holds that $\ex n{\N} p = \mbb p_n$.
\end{proof}

\begin{remark}[The use of \gls{LPO}]
  Notice that there is no way we can avoid the use of \gls{LPO} (or some other similar principle) in $\psh(\N)$ in the proof of~\zcref{prop:qcforFN}. From an external perspective, this does \emph{not} hold for $\psh(\N)$ in an arbitrary constructive meta-theory. By Kripke-Joyal semantics, it is easy to see that its validity depends on the fact that every upward closed subset of $\set{\ms p_0<\cdots<\ms p_n}$ is either empty or the principle upset $\dv \ms p_i$ for some $i\le n$, which only holds if we assume \gls{LEM} (or some other similar classical principle) in the meta-theory. The original work of~\citet{birkedal_first_2012} indeed uses classical reasoning in the meta-theory. It is interesting that in the work of~\citet{palombi_classifying_2023} a fully constructive development of the L\"ob induction for the topos of trees is developed. As we will see in~\zcref{subsec:lobindfortree}, the L\"ob induction in $\psh(\N)$ is indeed a consequence of~\zcref{prop:qcforFN}; however, this statement must be stronger, since it does not hold in an arbitrary constructive meta-theory. 
\end{remark}

\subsection{Elementary properties in the topos of trees}

In this section we work fully in the constructive internal logic of $\psh(\N)$. We assume it has a generic model $\mbb p = \set{\mbb p_n}_{n\in\N}$ of $\mbb F_\N$. We may also define $\ovob \coloneq \set{-1} \cup \N$, and define $\mbb p_{-1} \coloneq 0$. It turns our that for all our purpose, the following single axiom which corresponds to~\zcref{prop:qcforFN} suffices:

\PrintAxiomGR

We first see some other elementary consequences of (\AxiomGR) in $\psh(\N)$.

\begin{lemma}\label{lem:elementaryqcforFN}
  $\Omega$ is \emph{linear}, in the sense that 
  \[ \fa {p,q}\Omega (p \to q) \vee (q \to p). \]
  In particular, for any $n,m\in\ovob$,
  \[
  \mbb p_n \to \mbb p_m = 
  \begin{cases}
    1 & n \le m \\
    \mbb p_m & \other
  \end{cases} 
  \]
\end{lemma}
\begin{proof}
  Linearity is a trivial consequence of (\AxiomGR). The remaining holds by linearity and the fact that $\mbb p$ is a model of $\mbb F_\N$.
\end{proof}

\begin{remark}
  Notice that the second half of the statement in~\zcref{lem:elementaryqcforFN} can also be seen as a consequence of the \gls{QC} principle for a multi-sorted \emph{Horn theory} as in~\zcref{thm:qc}. Though the geometric theory $\mbb F_\N$ is not a Horn theory, it can indeed be presented as a \emph{quotient} of a Horn theory. Semantically, we could view $\psh(\N)$ as a subtopos of the presheaf topos $\psh(\N \cup \set{\infty})$. As the poset $\N \cup \set{\infty}$ has finite meets, $\psh(\N \cup \set{\infty})$ will indeed be the classifying topos of a multi-sorted Horn theory, and by~\zcref{thm:qcdescent} the corresponding \gls{QC} principle for this Horn theory will also hold in $\psh(\N)$, which governs the internal function spaces between definable sets.
\end{remark}

\begin{lemma}\label{lem:doublenegcontinuous}
  $\neg\neg \colon \Omega \to \Omega$ is equal to $p \mapsto 0 \to p$. In particular, it is continuous.
\end{lemma}
\begin{proof}
  By (\AxiomGR), it suffices to show this for all $\mbb p_n$ for $n\in\ovob$. When $n = -1$ this is trivial. When $n \ge 0$, by~\zcref{lem:elementaryqcforFN} we indeed have $\neg\neg\mbb p_n = \neg 0 = 1$, and $0 \to \mbb p_n = 1$. Hence, in both cases $\neg\neg$ coincides with $0 \to (-)$, and by function extensionality they coincides. 
\end{proof}

\begin{remark}\label{rem:empdecidable}
  From the computation in~\zcref{lem:doublenegcontinuous} we also see that $\Omega \cong 1 + \Omega_{\dneg}$, where $\Omega_{\dneg} \subseteq \Omega$ can be identified as the following set,
  \[ \Omega_{\dneg} = \scomp{p\in\Omega}{\ex n\N p = \mbb p_n}. \]
  The remaining truth value is simply $0 = \mbb p_{-1}$, and in particular, whether a proposition is false is \emph{decidable}.
\end{remark}

\begin{corollary}\label{Kuroda}
  The Kuroda principle holds: for any $p \colon X \to \Omega$, 
  \[ \fa xX \neg\neg p(x) \to \neg\neg \fa xX p(x). \]
\end{corollary}
\begin{proof}
  By~\zcref{lem:doublenegcontinuous}, $\neg\neg = 0 \to (-)$ which is a \emph{right adjoint}, hence preserves all limits.
\end{proof}

\begin{corollary}\label{negnegBoo}
  The double negation of \gls{LEM} holds: $\neg\neg\fa p{\Omega}p \vee \neg p$.
\end{corollary}
\begin{proof}
  By~\zcref{Kuroda}, it suffices to show for any $p\in\Omega$, $\neg\neg(p \vee \neg p)$, or equivalently, $\neg\neg p \vee \neg p$. This holds by linearity of $\Omega$ as shown in~\zcref{lem:elementaryqcforFN}.
\end{proof}

\begin{remark}\label{scatterlogic}
  \cite{ESAKIA200097} introduced the notion of \emph{$\bot$-scattered} and \emph{scattered} topos, and shows that the statements in \zcref{Kuroda} and~\zcref{negnegBoo} are in fact equivalent, in which case the topos is called $\bot$-scattered. As a consequence of~\zcref{canonicallob} and~\zcref{canonicalneg} later, we can in fact show internally that we are working with a \emph{scattered} topos.
\end{remark}

\begin{lemma}\label{internallyconnected}
  For any $n\in\N$, $\mbb p_n$ is connected, i.e.\ $2 \cong 2^{\mbb p_n}$. 
\end{lemma}
\begin{proof}
  Given any map $\mbb p_n \to 2$, this is equivalent to two propositions $p,q$ that $p \vee q = \mbb p_n$ and $p \wedge q = 0$. By (\AxiomGR), we may assume $p = \mbb p_i$ and $q = \mbb p_j$. This way, $\mbb p_i \wedge \mbb p_j = \mbb p_{i \wedge j} = 0$, which by~\zcref{rem:empdecidable}, $i \wedge j = -1$, equivalently $i = -1$ or $j = -1$. Thus, the map $\mbb p_n \to 2$ factors through some $1 \inj 2$, thus $2 \simeq 2^{\mbb p_n}$.
\end{proof}

In fact, we can show a stronger result:

\begin{lemma}\label{interallyconconnected}
  For any $n\in\omega$, $\N \cong \N^{\mbb p_n}$.
\end{lemma}
\begin{proof}
  Given $f \colon \mbb p_n \to \N$, again we can define a sequence of propositions $p \colon \N \to \Omega$, where for any $i\in\N$,
  \[ p(i) \coloneq f\inv(\set{i}). \]
  By construction, $\bigvee_{i\in\N}p(i) = \mbb p_n$ and $p(i) \wedge p(j) = 0$ for all $i \neq j$. The first condition implies that there exists $i\in\N$ that $p(i) = \mbb p_n$. The second condition implies that $p(j) = 0$ for all $j \neq i$. By~\zcref{rem:empdecidable}, such $i$ is unique, thus $f$ factors through a unique $\set i \hook \N$.
\end{proof}

\subsection{L\"ob induction in the topos of trees from quasi-coherence}\label{subsec:lobindfortree}

In this section we show that given (\AxiomGR), we can in fact \emph{define} an internal modality on $\Omega$ and prove it satisfies L\"ob induction. In fact, we will present two constructions. 

\begin{construction}[A positive construction of the later modality on propositions]\label{consprop}
  Let $\theta \colon \ovob \to \ovob$ be any \emph{monotone, inflationary} function: 
  \begin{itemize}
    \item for any $n,m\in\ovob$ with $n\le m$, $\theta n \le \theta m$;
    \item for any $n\in\ovob$, $n \le \theta n$.
  \end{itemize}
  Then the modality $\bx_\theta \colon \Omega \to \Omega$ takes $p\in\Omega$ to
  \[ \bx_\theta p \coloneq \bigvee_{n\in\ovob} \mbb p_{\theta n} \wedge (\mbb p_{n} \to p). \]
\end{construction}

\begin{proposition}\label{modalityonrepre}
  $\bx_\theta \colon \Omega \to \Omega$ makes the following diagram commute,
  \[ 
  \begin{tikzcd}
    \ovob \ar[d, two heads, "\mbb p_-"'] \ar[r, "\theta"] & \ovob \ar[d, two heads, "\mbb p_-"] \\
    \Omega \ar[r, "\bx_\theta"'] & \Omega
  \end{tikzcd}
  \]
\end{proposition}
\begin{proof}
  For any $n\in\ovob$, by~\zcref{consprop} we have 
  \begin{align*}
    &\bx_\theta\mbb p_{n} \\
    &\quad= \bigvee_{m\in\ovob} \mbb p_{\theta m} \wedge (\mbb p_{m} \to \mbb p_{n}) \\ 
    &\quad= \bigvee_{m \le n} \mbb p_{\theta m} \vee \bigvee_{m > n} \mbb p_{\theta m} \wedge \mbb p_{n} \\ 
    &\quad= \mbb p_{\theta n} \vee \mbb p_{n} \\ 
    &\quad= \mbb p_{\theta n}
  \end{align*}
  The first equality holds by definition; the second holds by~\zcref{lem:elementaryqcforFN}; the third holds by~\zcref{lem:elementaryqcforFN} and the fact that $\theta$ is monotone and inflationary; similarly for the last equality.
\end{proof}

\begin{remark}\label{usualmodalitypos}
  For the usual later modality $\bx$ in synthetic guarded domain theory as specified in~\citet{birkedal_first_2012}, it corresponds to the monotone inflation map $n \mapsto n\plo$ on $\ovob$. In particular, it can be internally defined as  
  \[ \bx p \coloneq \bigvee_{n\in\ovob}\mbb p_{n\plo} \wedge (\mbb p_{n} \to p). \]
  \zcref{lobprop} will provide an internal proof of L\"ob induction for this operator.
\end{remark}

As a start, we can show that the modality $\bx_\theta$ will always satisfy some desired property:

\begin{proposition}\label{posproplex}
  For any monotone and inflationary map $\theta \colon \ovob \to \ovob$, $\bx_\theta$ is pointed and left exact.
\end{proposition}
\begin{proof}
  To show $\fa p{\Omega}p \to \bx_\theta p$, by (\AxiomGR) it suffices to show $\mbb p_{n} \to \bx_\theta\mbb p_{n}$ for all $n\in\ovob$. By~\zcref{modalityonrepre} and~\zcref{lem:elementaryqcforFN}, this holds since $\bx_\theta\mbb p_{n} = \mbb p_{\theta n}$ and $\theta$ is inflationary. 
  
  For left exactness, by (\AxiomGR) we may assume $1 = \mbb p_{n}$, thus by~\zcref{modalityonrepre} we have $\bx_\theta 1 = \bx_\theta\mbb p_{n} = \mbb p_{\theta n}$. Since $\theta$ is inflationary, by~\zcref{lem:elementaryqcforFN} this means $\mbb p_{\theta n}$ holds as well, thus $\bx_\theta 1 = 1$. Proving $\bx_\theta$ preserves binary meets is similar.
\end{proof}

Most importantly, we have a criterion of when $\bx_\theta$ satisfies L\"ob induction:

\begin{theorem}\label{lobprop}
  If $\theta$ furthermore is strictly inflationary, viz.\ $n < \theta n$ for all $n\in\ovob$, then $\bx_\theta$ satisfies L\"ob induction:
  \[ \fa p{\Omega} (\bx_\theta p \to p) \to p. \]
\end{theorem}
\begin{proof}
  Take $p\in\Omega$, and by (\AxiomGR) let $p = \mbb p_{n}$. Assume $\bx_\theta p \to p$, viz. $\mbb p_{\theta n} \to \mbb p_{n}$. Since $\theta$ is strictly increasing, by~\zcref{lem:elementaryqcforFN} this implication is equivalent to $\mbb p_{n}$, viz.\ $p$.
\end{proof}

On the other hand, we can also provide a negative construction of the modality, which in fact corresponds more to the development in~\citet{ESAKIA200097}.

\begin{definition}[Flat map]\label{flat}
  We say a monotone function $\rho \colon \ovob \to \ovob$ is \emph{flat}, if it is unbounded, i.e.\ for any $n\in\ovob$, there exists $m\in\ovob$ that $n < \rho m$. 
\end{definition}

\begin{definition}\label{dualmapofflat}
  Given a flat map $\rho$, by the least number principle we have a dual map $\qsi\rho \colon \ovob \to \ovob$ that for any $n\in\ovob$, $\qsi\rho n$ is the least $m$ that $n < \rho m$. We also define the \emph{normalised} dual map $\pf\rho$ as $\pf\rho(n) \coloneq \max(\qsi\rho(n),n)$.
\end{definition}

\begin{remark}
  By definition, $\pf\rho$ is monotone and inflationary.
\end{remark}

\begin{construction}[A negative construction of the later modality on propositions]\label{consnegprop}
  Let $\rho \colon \ovob \to \ovob$ be a \emph{flat} map.
  Then the modality $\bx^\rho \colon \Omega \to \Omega$ takes $p:\Omega$ to
  \[ \bx^\rho p \coloneq \bigwedge_{n:\ovob} \mbb p_{n} \vee (\mbb p_{\rho n} \to p). \]
\end{construction}

\begin{remark}
  The form of the construction of modality in~\zcref{consnegprop} is evidently \emph{dual} to that in~\zcref{consprop}.
\end{remark}

Again, let us compute the value of $\bx^\rho$ on the generic propositions.

\begin{proposition}\label{neglateronrep}
  Let $\rho$ be a flat map and $\pf\rho$ its normalised dual map. For any $n\in\ovob$,
  \[ \bx^\rho\mbb p_{n} = \mbb p_{\pf\rho n}. \]
\end{proposition}
\begin{proof}
  By~\zcref{lem:elementaryqcforFN},~\zcref{consnegprop} and the definition of $\qsi\rho,\pf\rho$ in~\zcref{dualmapofflat},
  \begin{align*}
    \bx^\rho\mbb p_{n} 
    &= \bigwedge_{m\in\ovob} \mbb p_{m} \vee (\mbb p_{\rho m} \to \mbb p_{n}) \\ 
    &= \bigwedge_{m \ge \qsi\rho n} \mbb p_{m} \vee \mbb p_{n} \\
    &= \mbb p_{\qsi\rho n} \vee \mbb p_{n}  \\
    &= \mbb p_{\pf\rho n} \qedhere
  \end{align*}
  The first identity holds by definition. The second holds since when $m < \qsi\rho n$, by definition $\rho m \le n$ thus the expression evaluates to 1; and by (QC), when $\rho m > n$, $\mbb p_{\rho m} \to \mbb p_{n}$ is equivalent to $\mbb p_{n}$. The last identity holds by definition of $\pf\rho$.
\end{proof}

\begin{corollary}\label{negispos}
  For any flat map $\rho$, $\bx^\rho = \bx_{\pf\rho}$.
\end{corollary}
\begin{proof}
  By function extensionality, it suffices to show $\bx^\rho p = \bx_{\pf\rho}p$ for all $p:\Omega$. Now by (\AxiomGR) it suffices to show this for all $\mbb p_{n}$. The remaining holds by~\zcref{neglateronrep} and~\zcref{modalityonrepre}.
\end{proof}

\begin{corollary}
  For any flat map $\rho$, $\bx^\rho$ is pointed and left exact.
\end{corollary}

\begin{theorem}\label{canonicallob}
  If $\rho$ is deflationary, i.e. $\rho n \le n$ for all $n:\ovob$, then $\bx^\rho$ satisfies L\"ob induction.
\end{theorem}
\begin{proof}
  If $\rho$ is deflationary, then by monotonicity $\pf\rho$ must be strictly inflationary. Thus the claim holds by~\zcref{lobprop} and~\zcref{negispos}.
\end{proof}

\begin{remark}\label{canonicalneg}
  By~\zcref{usualmodalitypos} and~\zcref{negispos}, we have an alternative construction of the usual later modality $\bx$. Notice that the identity map $\id \colon \ovob\to\ovob$ is flat, and its normalised dual map is exactly given by $n \mapsto n\plo$. This way, we have an alternative characterisation of $\bx$ as below,
  \[ \bx p \coloneq \bigwedge_{n\in\ovob} \mbb p_{n} \vee (\mbb p_{n} \to p). \]
  This nice thing about this formula is that, by (\AxiomGR), we can avoid mentioning the generic propositions at all,
  \[ \bx p = \bigwedge_{q\in\Omega} q \vee (q \to p). \]
  This is indeed the formula obtained in~\cite{ESAKIA200097}.
\end{remark}

\begin{corollary}
  $\psh(\N)$ is a scattered topos.
\end{corollary}
\begin{proof}
  This follows from~\citet[Thm.~2]{ESAKIA200097} the fact that the modality $\bx p = \bigwedge_{q\in\Omega} q \vee (q \to p)$ satisfies L\"ob induction from~\zcref{canonicallob} as mentioned in~\zcref{canonicalneg}.
\end{proof}

\begin{remark}[Modality on sets]
  It is also possible to construct the later modality on all \emph{sets} and show it admits the usual guarded fixed point construction purely from the internal logic with (\AxiomGR). Roughly speaking, we may define the modality on sets by an end or 
\end{remark}

\section{Application II: quantifier elimination and Boolean classifying topoi}

\begin{definition}[Geometric theories admiting ]
  
\end{definition}

\section{Outlooks}

How should we imagine a quasi-coherence principle for the \emph{set universe}?